\section{Resultados}
\label{sec:resultados}

Esta seção será preenchida após a execução dos experimentos e coleta das métricas. Nesta versão inicial, as tabelas indicam quais informações serão consolidadas no artigo final.

\begin{table}[ht]
\centering
\caption{Caracterização do ambiente experimental}
\label{tab:ambiente}
\begin{tabular}{ll}
\toprule
\textbf{Item} & \textbf{Descrição} \\
\midrule
CPU & A preencher \\
Memória RAM & A preencher \\
GPU & A preencher \\
Armazenamento & A preencher \\
Sistema operacional & A preencher \\
Ferramentas de medição & A preencher \\
Carga experimental & A preencher \\
\bottomrule
\end{tabular}
\end{table}

\begin{table}[ht]
\centering
\caption{Métricas de execução previstas para análise}
\label{tab:metricas}
\begin{tabular}{lll}
\toprule
\textbf{Execução} & \textbf{Métrica} & \textbf{Valor observado} \\
\midrule
Execução 1 & Tempo total & A preencher \\
Execução 1 & Uso de CPU & A preencher \\
Execução 1 & Uso de memória & A preencher \\
Execução 1 & Uso de GPU & A preencher \\
Execução 1 & Observação principal & A preencher \\
\bottomrule
\end{tabular}
\end{table}

Após o preenchimento das tabelas, os resultados serão analisados buscando responder se a carga executada foi limitada principalmente por processamento, memória, armazenamento, aceleração ou outro fator observado durante os testes.
